\scanpage{28}
in $\mathcal E$, which one can choose canonically by taking the quarter of a great circle which joins $P$ to the midpoint $2$ of the \deleted{segment} the arc $(1,\infty)$ of a great circle opposite $0$, which defines an isomorphism of the universal covering of $U$ at $P$, with that of $U$ at $\bar P$ on which the preceding group acts in a canonical way, hence also by transport of structure on $\tilde U\deleted{\unsure{(\overline P)|(\mathcal E)}}_{\deleted{\unsure{\mathcal E|\xi}}}(P)$, hence it maps into $\mathcal E$.

One finds that
\[
\dot\sigma_0\mapsto\sigma_0,\quad\dot\tau\to\tau_0
\]
is the corresponding lifting. Likewise for the fixed point $-1$ of $\dot\sigma_1$, or $\frac12$ of $\dot\sigma_\infty$, one finds partial splittings over the \deleted{\unsure{conjugacy}} stabilizer subgroups $\Gamma_{-1}\subset\Gamma$ and $\Gamma_2\subset\Gamma$. It is convenient to put
\[
\deleted{\unsure{\Gamma|F}}\quad 2=0',\quad -1=1',\quad \tfrac12=\infty',
\]
whence isotropy groups $\Gamma_{0'}=\{\dot\sigma_0,\tau\}$, $\Gamma_{1'}=\{\dot\sigma_1,\tau\}$, and $\Gamma_{\infty'}=\{\dot\sigma_\infty,\tau\}$, which lift resp to $\sigma_0$ and $\tau_0$, \unsure{to|or} $\sigma_1$ and $\tau_1$, to $\sigma_\infty$ and $\tau_\infty$.

The connected components of $U^\tau$ are the three great-circle arcs between $0,1,\infty$, which each contain one \deleted{point} and only one of the points $0',1',\infty'$, and the conjugacy classes
\mspage{27}
